% Target language: Korean (ko-KR).
% Translation source: American English, commit 5da89be81487b8164a2f541c7245ba75ba07be53.
% Checked against Simplified Chinese, commit 0d064f6ce5c6f02a2e89edc6004b56723e8042be.
% Status: AI-assisted translation; needs native-speaker review.
\StoryParagraph{3}{난자 하나의 무게, 3.5마이크로그램.}

\StoryParagraph{3}{현금 5천만 달러의 무게, 2.5톤.}

\StoryParagraph{3}{몽타주 라구나 비치에서 걸려 온 전화로 징톈이 내게 요구한 것은 후자였고, 담보로 잡은 것은 전자였다.}

\StoryParagraph{3}{그녀가 그곳에 간 건 난자 채취를 위해서였다. 대리모가 낳을 아이는 2027년에 태어날 예정이었다. 서두르면 아직 말띠로 태어날 수도 있었다. 그 이야기를 먼저 꺼낸 것도 그녀였다. 병원에 도착하자 그녀는 5천만 없이는 채취하지 않겠다고 했다.}

\UnprintedPause

\StoryParagraph{0}{몽타주 라구나 비치는 태평양을 내려다보는 해안 절벽 위에 있다. 우리는 그녀의 사생활을 위해 한 층 전체를 빌렸다. 30일, 100만 달러. 그녀는 8일 만에 떠났다.}

\StoryParagraph{0}{남은 22일 동안에도 호텔 청소부들은 날마다 평소처럼 그 층에 올라갔다. 비서에게 들은 이야기다. 열쇠를 돌려주러 갔다가 한 번 본 적이 있다고 했다. 멕시코인 여자가 카트를 밀며 방마다 들어가 수건을 갈고, 상판을 닦고, 시트 모서리를 여미고, 뒤로 물러나와 문을 닫았다. 한 층을 전부 끝내는 데 두 시간.}

\StoryParagraph{0}{비서는 아무도 묵지 않는데 왜 계속 청소하느냐고 물었다.}

\StoryParagraph{0}{그 여자는 투숙객 명부에 누군가 묵고 있는 것으로 적혀 있다고 했다.}

\StoryParagraph{0}{왕비를 잃은 베르사유 궁전.}

\StoryParagraph{0}{그때 문득 알았다. 징톈은 이런 것을 좋아했다. 자신 때문에 어떤 장소가 텅 빈 채로 있는 것을. 아무 일도 일어나지 않는 것을.}

\StoryParagraph{3}{몽타주 라구나 비치와 함께 비어 있던 것은 밴나이즈에 멋쩍게 서 있는 걸프스트림 G700이었다. 승무원들은 날마다 평소처럼 출근했다.}

\StoryParagraph{3}{8일째에 그녀가 그 전화를 걸었다. 그 뒤로 그녀를 본 사람은 아무도 없었다.}

\NarrativeBreak

\StoryParagraph{0}{2007년에는 이미 누군가 그녀를 위해 돈을 태우고 있었다. 내가 베이징대학교에 막 입학한 해였다.}

\StoryParagraph{3}{물론 돈을 태운 사람은 내가 아니었다.}

\StoryParagraph{3}{QQ에 팝업 창이 떴고, 광고 모델은 그녀였다. 낡은 Lenovo IBM 컴퓨터였고 해상도는 흐릿했다. 나는 이미지를 저장했다.}

\StoryParagraph{3}{1년 기숙사비가 1천 위안이던 베이징대학교 기숙사에서 인터넷에 접속해 샤오네이에서 그녀를 찾아 친구 신청을 보냈다. 샤오네이는 메시지를 한마디 덧붙이게 했다. 쓰고 지우고, 지우고 다시 쓰기를 40분쯤 했다. 마지막에 보낸 말은 “안녕하세요.”였다.}

\StoryParagraph{3}{그녀는 받아 주지 않았다.}

\StoryParagraph{3}{어쩌면 샤오네이를 한 번도 쓰지 않았을지도 모른다.}

\StoryParagraph{3}{그 뒤 여러 해 동안 나는 무슨 일을 하든 한 가지를 생각했다. 어디까지 가야 그녀가 받아 줄까. 첫 10억 때도 생각했다. 100억 때도 생각했다.}

\StoryParagraph{3}{나중에는 생각하지 않았다. 런런왕은 진작 망해 있었다.}

\StoryParagraph{3}{두 번째 신청은 보내지 않았다.}

\StoryParagraph{3}{2011년, 베이징은 추웠다. 나는 150위안짜리 패딩을 샀다. 사흘을 망설였기에 그 가격을 기억한다.}

\StoryParagraph{3}{같은 해 그녀는 제작비 1억 5천만 위안의 영화 《전국: 천하영웅의 시대》에 출연했다. 쑨훙레이, 우전위, 김희선, 나카이 키이치, 강무는 19년 뒤의 몽타주 라구나 비치와 걸프스트림 G700처럼 그녀를 돋보이게 하려고 불려 왔다. 그들은 자신이 왜 거기 있는지 몰랐다. 멀리까지 왔으니 무슨 일인가 일어나야 했다. 아무 일도 일어나지 않았다.}

\StoryParagraph{3}{그다음에는 1억 5천만 달러짜리 《그레이트 월》이 있었다. 장이머우. 유덕화. 맷 데이먼. 그다음은 성룡. 징톈 우주의 명단은 점점 길어졌다.}

\StoryParagraph{3}{그 무렵 나는 이미 미국에 있었다. 공부하고, 창업하고, 갈수록 바빠졌다. 한 가지만은 변하지 않았다. 기기를 바꿀 때마다 아무리 번거로워도 그 QQ 이미지는 반드시 옮겼다.}

\StoryParagraph{3}{19년 뒤, 그 명단에 내 이름이 더해졌다.}

\NarrativeBreak

\StoryParagraph{3}{우리는 로즈우드 홍콩에서 처음 만났다. 그녀를 만나는 날이면 긴장할 줄 알았다.}

\StoryParagraph{3}{긴장하지 않았다.}

\StoryParagraph{3}{그녀는 사진보다 말랐고, 편한 옷차림으로 소파에 기대어 있었다. 그녀를 보고 있자니 내가 정확히 무엇을 기다려 왔는지 문득 기억나지 않았다.}

\StoryParagraph{3}{나는 그녀를 카텔란의 바나나를 보러 데려갔다. 그녀는 한참 바라보다가 물었다. “이거 얼마야?”}

\StoryParagraph{3}{620만 달러.}

\StoryParagraph{3}{그녀가 말했다. “썩으면 어떻게 해?”}

\StoryParagraph{3}{나는 말이 없었다. 그녀가 말했다. “됐어. 영화나 보러 가자.”}

\StoryParagraph{3}{나는 “그래.”라고 했다.}

\StoryParagraph{3}{그녀는 상영관에 우리 둘 말고는 아무도 있어서는 안 된다고 했다.}

\StoryParagraph{3}{나는 “그래.”라고 했다.}

\StoryParagraph{3}{나는 전화를 걸러 밖으로 나갔다. 돌아왔을 때도 그녀는 그 바나나 앞에 서 있었다.}

\StoryParagraph{3}{영화는 《주토피아 2》였다.}

\StoryParagraph{3}{두 사람만을 위한 상영관. 이런 광경은 보통 중앙선전부가 영화를 심사할 때나 쓰인다.}

\StoryParagraph{3}{나는 비서에게 영화관의 모든 좌석을 사게 했다.}

\StoryParagraph{3}{가짜 흥행 성적을 비판하는 글을 읽은 적이 있다. 표는 매진인데 관객은 두 명뿐이라는 이야기였다. 이제는 알 것 같았다. 언제나 흥행 성적 때문인 것은 아니었다.}

\StoryParagraph{3}{좌석을 전부 사는 일은 어렵지 않았다. 문제는 이미 팔린 표였다. 비서들과 경호원들은 일찍 극장 입구에 가서 기다렸다. 표를 가진 관객이 나타날 때마다 현금을 건넸다.}

\StoryParagraph{3}{아이를 데려온 부부가 왔다. 아이는 벌써 입구에 서서 팝콘 통을 높이 들고 있었다.}

\StoryParagraph{3}{비서가 돈뭉치를 내밀었다. 아이는 왜 들어갈 수 없는지 이해하지 못했다. 부모는 돈을 받고, 고맙다고 말하고, 아무것도 묻지 않은 채 떠났다.}

\StoryParagraph{3}{그날 돈을 받은 사람은 26명이었다.}

\StoryParagraph{3}{영화를 본 뒤 그녀는 위챗 프로필 사진을 주디 홉스 경관으로 바꿨다.}

\StoryParagraph{3}{그날 밤 나는 손으로 그녀의 가슴을 주물렀다. 그녀가 나를 밀쳐 냈다.}

\StoryParagraph{3}{싫은 줄 알았다. 손을 거두고 무슨 말인가 하려 했다.}

\StoryParagraph{3}{그녀는 아무 말도 하지 않았다. 몸을 일으켜 가방에서 작은 주머니를 찾아 지퍼를 열었다. 안에는 아주 작은 공구 세트 같은 것들이 한 줄로 놓여 있었다.}

\StoryParagraph{3}{그녀가 말했다. “네 손톱에 찔려서 아파.”}

\StoryParagraph{3}{내 손톱을 보았다. 나는 늘 손톱깎이로 되는대로 자를 뿐이어서 가장자리가 날카로웠다. 평소에는 느끼지 못했다.}

\StoryParagraph{3}{그녀가 말했다. “이리 와.”}

\StoryParagraph{3}{그녀는 내 손을 자기 무릎 위에 놓고 손톱을 하나씩 갈았다.}

\StoryParagraph{3}{먼저 거친 면. 다 갈면 다른 면으로, 다시 또 다른 면으로 바꿨다. 한동안 갈고는 손가락을 불빛에 비춰 보고, 다시 갈았다.}

\StoryParagraph{3}{나는 얼마나 오래 갈아야 하느냐고 물었다.}

\StoryParagraph{3}{그녀가 말했다. “사람을 긁지 않을 때까지.”}

\StoryParagraph{3}{손가락 열 개. 그녀는 거의 한 시간을 들였다.}

\StoryParagraph{3}{그러는 동안 그녀가 가르쳐 주었다. 이 면으로 모양을 잡고, 이 면으로 윤을 낸다. 마지막 면은 앞뒤로 문지르지 말고 한 방향으로 움직여야 한다.}

\StoryParagraph{3}{나는 “알았어.”라고 했다.}

\StoryParagraph{3}{다 끝나자 그녀는 내 손을 뒤집어 살펴보고 손끝으로 쓸어 보았다.}

\StoryParagraph{3}{그녀가 말했다. “됐어.”}

\StoryParagraph{3}{“만져. 어디든 네가 만지고 싶은 데.”}

\StoryParagraph{3}{도중에 그녀가 한마디 했다. “언젠가 내가 없어지면, 다시는 아무도 네 손톱을 다듬어 주지 않을 거야.”}

\StoryParagraph{3}{나는 살짝 웃고 아무 말도 하지 않았다.}

\StoryParagraph{0}{그녀를 만난 지 석 달쯤 되었을 때, 나는 사람을 시켜 자료를 하나 만들었다. 그녀가 그동안 해 온 사업, 협력자, 회사 지분, 징궈칭, 톈신아이, 관련된 모든 계좌, 베이징에서 빌린 집, 상하이에서 산 집, 그 밖의 몇 가지. 일주일 만에 완성됐고 모두 40여 쪽이었다. 남녀 관계에 관한 일도 있었다. 내게는 그것이 가장 중요하지 않았다.}

\StoryParagraph{0}{그 안의 몇몇 일은 나중에 그녀가 직접 말해 주었다. 그녀가 한 말은 거기에 적힌 것과 조금 달랐다. 나는 아무 말도 하지 않았다.}

\StoryParagraph{0}{끝내 말하지 않은 일도 있었다. 나도 묻지 않았다.}

\StoryParagraph{0}{나는 그녀가 하는 말을 전부 믿었다.}

\StoryParagraph{0}{그 자료는 간직했다. 지금도 가지고 있다.}

\StoryParagraph{3}{하지만 한 글자도 믿지 않았다.}

\NarrativeBreak

\StoryParagraph{3}{그녀가 말했다. “홍콩은 사람이 너무 많아. 알아보는 사람이 있을 거야. 어디 좀 가자.”}

\StoryParagraph{3}{나는 “그래.”라고 했다.}

\StoryParagraph{3}{그녀는 지도에서 테네리페를 가리켰다.}

\StoryParagraph{3}{나는 “그래.”라고 했다.}

\StoryParagraph{3}{전화를 걸러 나갔다가 돌아와서야 그곳이 아프리카 북서부 해안에서 떨어진 대서양에 있다는 걸 알았다. 다음 날 정오, 우리 일행 30명은 출발했다. 공항에는 A330이 기다리고 있었다.}

\StoryParagraph{3}{비행기 침실에서 그녀가 물었다. “내 이름, 예뻐?”}

\StoryParagraph{3}{나는 예쁘다고 했다.}

\StoryParagraph{3}{그녀가 말했다. “아빠 성과 엄마 성을 붙인 거야. 참 신기하지?”}

\StoryParagraph{3}{나는 첫날부터 기억했다고 했다.}

\StoryParagraph{3}{그녀는 나를 안고 말했다. “앞으로 징톈이라고 부르지 말고 엄마라고 불러.”}

\StoryParagraph{3}{“알았어, 엄마.”}

\StoryParagraph{3}{착륙한 뒤 기장이 다가와 말했다. “사장님, 처음에는 섬 쪽에서 착륙을 허가하지 않았습니다. 이렇게 큰 전용기를 본 적이 없답니다.”}

\StoryParagraph{3}{섬에서도 규칙은 같았다. 우리 주변에 사람이 있어서는 안 됐다. 비서들이 처리했다. 아무도 이유를 묻지 않았다. 이번에는 나도 몇 명이었는지 묻지 않았다.}

\StoryParagraph{3}{리츠칼튼은 대서양을 마주 보고 있었고 통유리창은 너무 컸다. 그녀는 한 줄기 빛만 있어도 잠들지 못했다. 비서들은 방 전체를 검은 테이프로 덮었다. 밖은 대낮처럼 밝았고 안은 독방처럼 어두웠다. 테이프는 햇볕을 받으면 떨어져서 날마다 다시 붙여야 했다.}

\StoryParagraph{3}{비서는 모두 11명이었다. 아침 6시 30분에 와서 우리가 잠든 뒤에야 떠났다. 테이프를 붙이는 사람은 여자 둘이었다. 한 명은 창틀을, 다른 한 명은 에어컨과 콘센트의 표시등을 맡았다. 다 붙인 뒤에는 세 번씩 확인했다. 테이프를 30롤 샀다가 나중에 20롤을 더 샀다. 오후 3시가 지나 붙인 테이프가 가장 쉽게 떨어진다는 걸 그들이 알아냈다. 그때 유리가 가장 뜨거워졌기 때문이다. 그래서 아침에 붙이기로 했다. 아무도 가르쳐 주지 않았다. 직접 해 보며 알아낸 일이었다. 그들은 마치 하나의 공정을 논하듯 서로 의논했다.}

\StoryParagraph{3}{그녀는 그 두 여자의 이름을 알았다. 나는 몰랐다.}

\StoryParagraph{3}{우리는 테이데 국립공원 길가에 차를 세우고 해가 지는 모습을 보았다. 비서들은 차에서 장미를 내리고 샴페인을 열었다.}

\StoryParagraph{3}{엄마는 구름 바다를 바라보며 말했다. “정말 예쁘다. 그래도 칭하이성 더링하의 노을만큼은 아니야.”}

\StoryParagraph{3}{나는 칭하이 사람이다.}

\StoryParagraph{3}{말하지 않았다.}

\StoryParagraph{3}{밤에는 산꼭대기에서 별을 보았다. 두툼한 솜옷을 입고도 추웠다.}

\StoryParagraph{3}{엄마가 말했다. “소원 빌어.”}

\StoryParagraph{3}{나는 무슨 소원을 빌었느냐고 물었다.}

\StoryParagraph{3}{엄마가 말했다. “안 알려 줄 거야.”}

\StoryParagraph{3}{엄마가 말했다. “너무 예쁘다. 그런데 너무 추워. 우리 몰디브에 갈까?”}

\StoryParagraph{3}{나는 “그래.”라고 했다.}

\StoryParagraph{3}{테네리페는 영하 15도, 몰디브는 영상 27도였다. 직선거리 9,500킬로미터, 비행시간 11시간. 하지만 몰디브 공항이 너무 작아서 가는 도중 광동체기에서 작은 비행기로 갈아타야 했다.}

\StoryParagraph{3}{엄마가 말했다. “작은 비행기는 너무 시끄러워. 큰 비행기가 좋아.”}

\StoryParagraph{3}{몰디브로 가는 비행기의 침실에서 엄마는 소원이 나와 아이를 갖는 것이라고 말했다. 반드시 대리 출산이어야 했다.}

\StoryParagraph{3}{나는 이유를 묻지 않았다. 아픈 게 두려운 줄 알았다.}

\StoryParagraph{3}{비행 도중 엄마가 갑자기 일어나 휴대전화를 찾았다. 스타링크에 연결하더니 말했다. “큰일 났어.”}

\StoryParagraph{3}{나는 무슨 일이냐고 물었다.}

\StoryParagraph{3}{그녀가 말했다. “나를 협박하는 사람이 있어. 사흘째 연락하는데 계속 답을 안 했어.”}

\StoryParagraph{3}{나는 얼마를 원하느냐고 물었다.}

\StoryParagraph{3}{그녀가 말했다. “넌 상관하지 마.”}

\StoryParagraph{3}{그녀는 오랫동안 글자를 입력했다. 지우고 고치고, 고치고 지웠다. 다 쓰고 나서는 휴대전화를 뒤집어 허벅지 위에 놓았다. 잠시 뒤 다시 돌려서 들여다보았다.}

\StoryParagraph{3}{나는 그 사람이 누구냐고 물었다.}

\StoryParagraph{0}{그녀가 말했다. “전에 나와 일하던 사람이야. 재무 담당. 감옥에 갔었어. 대신 죄를 뒤집어썼지.”}

\StoryParagraph{3}{“나온 뒤에는 갈 데가 없어서, 기어이 나를 찾아오겠대.”}

\StoryParagraph{3}{나는 차라리 경찰에 신고해서 잡아가게 하라고 했다.}

\StoryParagraph{3}{그녀가 나를 보았다. “넌 몰라.”}

\StoryParagraph{3}{나는 돈을 줬느냐고 물었다.}

\StoryParagraph{3}{그녀는 줬다고 했다.}

\StoryParagraph{3}{얼마.}

\StoryParagraph{3}{8만.}

\StoryParagraph{3}{나는 그걸 사흘이나 고민했느냐고 물었다.}

\StoryParagraph{3}{그녀는 말이 없었다.}

\StoryParagraph{3}{그날 남은 시간 동안 그녀는 거의 말하지 않았다. 착륙하고 나서야 잠들었다.}

\StoryParagraph{0}{소네바 자니에 가장 먼저 도착한 것은 엄마의 여행 가방이었다. 30명의 짐은 작은 비행기 한 대에 다 실리지 않아 세 대에 나눠 실었다. 엄마에게는 여행 가방 하나뿐이었다. 그 가방 하나 때문에 비행기 한 대를 따로 띄웠다. 다른 사람의 짐과 함께 두지 말라고 했기 때문이다.}

\StoryParagraph{0}{소네바 자니는 “no news, no shoes”로 유명했다. 엄마는 안심했다.}

\StoryParagraph{3}{호텔 버틀러, 경호원, 비서를 합해 26명이 붙었다. 드론 조종사 두 명이 멀찍이 따라왔다. 비서들은 다시 테이프를 붙이기 시작했다. 에어컨 불빛, 콘센트 불빛. 몇 번이고 되풀이했다.}

\StoryParagraph{0}{밤이면 우리는 암호를 주고받듯 서로의 옛 연인 이야기를 주고받았다. 마치 한 왕조의 연호를 하나씩 맞춰 보는 듯했다. 엄마는 어떤 어긋남도 허락하지 않았다. 다만 엄마에게는 금기어가 하나 있었다. 그 운동선수 이야기는 꺼내면 안 됐다.}

\StoryParagraph{0}{나는 꺼내지 않았다. 엄마의 등에 있는 문신이 그 운동선수 때문에 새긴 것임을 알면서도.}

\StoryParagraph{3}{너무 커서 지우지 못했는지도 모른다.}

\StoryParagraph{0}{그녀가 말했다. “이제 우리는 가족이잖아. 우리 부모님이 나를 이만큼 키우느라 고생하셨어. 부모님께 혼인 예물금을 드려야 해.”}

\StoryParagraph{0}{나는 “알았어.”라고 했다.}

\StoryParagraph{0}{그녀는 징궈칭과 톈신아이라는 두 이름과 두 계좌번호를 불러 주었다.}

\StoryParagraph{0}{3천만. 인민폐로.}

\StoryParagraph{0}{나는 송금했다. 그 두 이름은 40여 쪽짜리 자료에서 처음 보았다.}

\StoryParagraph{3}{잠시 뒤 그녀가 휴대전화를 내밀었다. 톈신아이의 답장은 한마디였다. “받았습니다.”}

\StoryParagraph{0}{한번은 엄마가 말했다. “내가 이런 걸 요구할 때 무슨 생각을 하는지 알아?”}

\StoryParagraph{0}{나는 무슨 생각이냐고 물었다.}

\StoryParagraph{0}{그녀가 말했다. “언젠가는 네가 싫다고 말할까 생각해.”}

\StoryParagraph{0}{나는 그러지 않을 거라고 했다.}

\StoryParagraph{0}{그녀가 말했다. “그럼 재미없어.”}

\StoryParagraph{3}{내가 못 할 일이 대체 무엇이 있을까, 하고 생각했다.}

\StoryParagraph{0}{너무 심심할 때면 나는 2007년의 QQ 이미지를 만지작거리며 새 기기로 옮기려 했다. 엄마가 얼굴을 들이밀었다. “왜 나 몰래 다른 여자 사진을 간직하고 있어? 누구야?”}

\StoryParagraph{0}{나는 침묵했다. 아무 말도 하지 않았다.}

\StoryParagraph{0}{사진이 너무 흐려서 그녀는 다른 일을 하러 갔다.}

\StoryParagraph{0}{그날 밤 나는 그 이미지를 지웠다.}

\StoryParagraph{3}{다음 날 오래된 휴대전화에서 다시 찾아냈다.}

\StoryParagraph{0}{엄마도 계속 휴대전화를 넘겨 보고 있었다. 나는 무엇을 보느냐고 물었다.}

\StoryParagraph{0}{그녀가 말했다. “아무것도 아니야. 주연 서열을 두고 싸우는 그 지겨운 일들.”}

\StoryParagraph{0}{주연 서열 다툼. 한 작품에서 맨 앞에 이름을 올릴 수 있는 스타는 한 명뿐이다. 남녀 주연 사이에서도 먼저 이름을 올릴 수 있는 사람은 한 명뿐이다.}

\StoryParagraph{0}{나는 지루하다고 했다.}

\StoryParagraph{0}{그녀가 말했다. “그러게. 더 힘든 건 가끔 촬영장에서 부부처럼 지낸다는 거야. 촬영에 들어가면 40일, 50일 동안 대개 남자 한 명과 연기하잖아. 밤에는 섹스하고 낮에는 누가 먼저 이름을 올릴지 싸워.”}

\StoryParagraph{0}{나는 인간 본성에 너무 어긋난다고 했다.}

\StoryParagraph{0}{그녀가 그렇다고 했다.}

\StoryParagraph{0}{“하지만 촬영이 끝나면, 그전에 몇 번을 섹스했든 아무 일도 없었던 것처럼 굴어.”}

\StoryParagraph{3}{그러고도 그녀는 30분을 더 읽었다.}

\StoryParagraph{0}{2026년 1월 2일, 소네바 자니. 새해에 나는 엄마에게 청혼했다.}

\StoryParagraph{0}{결심한 것은 1월 1일 아침이었다. 반지는 홍콩에 있었다. 섬에는 반지를 파는 곳이 없어서 비서가 수상비행기를 타고 말레로 갔다. 오후에 돌아왔고, 반지는 호텔 봉투 안에 들어 있었다.}

\StoryParagraph{0}{엄마가 말했다. “다이아몬드가 너무 작아.”}

\StoryParagraph{0}{엄마가 말했다. “괜찮아. 나중에 더 큰 걸로 다시 해 주면 돼.”}

\StoryParagraph{0}{“그리고, 베이징에 올 수 있어?”}

\StoryParagraph{0}{나는 침묵했다.}

\StoryParagraph{0}{그녀는 내가 가기 싫은 줄 알았다.}

\StoryParagraph{0}{엄마가 말했다. “왜 나한테 이렇게 잘해 줘?”}

\StoryParagraph{0}{나는 좋아하니까 그렇다고 했다.}

\StoryParagraph{0}{그녀가 말했다. “전에도 다른 사람한테 이랬어?”}

\StoryParagraph{0}{나는 아니라고 했다.}

\StoryParagraph{3}{그녀는 살짝 웃고 아무 말도 하지 않았다.}

\StoryParagraph{0}{어느 날 아침 일찍 눈을 떴다. 그녀는 침대에 없었다.}

\StoryParagraph{0}{밖으로 나가 보니 그녀는 테라스에 앉아 바다를 바라보고 있었다. 바다 전체가 검었다. 불은 켜지지 않았고 아직 동이 트지 않았다.}

\StoryParagraph{0}{나는 왜 이렇게 일찍 일어났느냐고 물었다.}

\StoryParagraph{0}{그녀가 말했다. “날짜를 세고 있어.”}

\StoryParagraph{0}{나는 무엇을 세느냐고 물었다.}

\StoryParagraph{0}{그녀가 말했다. “돌아가서 자.”}

\StoryParagraph{0}{나는 그녀 곁에 잠시 앉아 있었다. 무슨 말을 해야 할지 몰랐다. 그러다 돌아가서 잤다.}

\StoryParagraph{3}{다시 눈을 떴을 때 그녀는 이미 화장을 마친 뒤였다.}

\StoryParagraph{0}{몰디브에서 14일을 지내는 동안 그녀를 알아본 사람은 한 명도 없었다.}

\StoryParagraph{0}{13일째 밤, 식당에서 그녀는 선글라스를 벗었다. 잠시 뒤에는 모자도 벗어 탁자 위에 놓았다. 식사 내내 그대로 앉아 있었다. 옆 탁자에는 독일인 노부부가 있었고, 종업원은 스리랑카 사람이었다. 아무도 그녀를 다시 쳐다보지 않았다.}

\StoryParagraph{0}{식사를 마치자 그녀는 모자를 다시 썼다.}

\StoryParagraph{0}{엄마가 말했다. “돌아가자.”}

\StoryParagraph{0}{나는 아직 더 있기로 하지 않았느냐고 물었다.}

\StoryParagraph{3}{그녀가 말했다. “너무 조용해.”}

\UnprintedPause

\StoryParagraph{0}{홍콩으로 돌아오자 엄마는 내 부모님을 만났다.}

\StoryParagraph{0}{그날 밤 그녀는 양가 상견례와 결혼 준비를 해야 하니 가능한 한 많이 보내 달라고 했다. 여전히 징궈칭과 톈신아이였다.}

\StoryParagraph{0}{나는 이유를 묻지 않았다.}

\StoryParagraph{0}{나는 카드 한 장당 하루 한도가 100만이라고 했다.}

\StoryParagraph{0}{엄마가 말했다. “너무 느려.”}

\StoryParagraph{0}{“다섯 장 써.”}

\StoryParagraph{0}{카드 다섯 장, 계좌 두 개. 번갈아 송금했다. 세 번째 카드에서 나는 잠시 망설였다. 그래도 끝까지 모두 보냈다.}

\StoryParagraph{3}{엄마가 말했다. “너무 비효율적이야.”}

\StoryParagraph{0}{나는 춘제에는 양쪽 부모님을 모두 홍콩으로 모셔 와 함께 보내자고 했다.}

\StoryParagraph{0}{엄마가 말했다. “올해는 이미 하이난에서 보내기로 했어. 내년에 하자.”}

\StoryParagraph{0}{춘제에 빅토리아 하버에서 불꽃놀이가 열렸고 100만 명이 모였다. 그녀는 없었다.}

\StoryParagraph{0}{새벽 3시, 나는 부모님을 집에 모셔다드렸다.}

\StoryParagraph{0}{길은 텅 비어 있었다. 빨간불 앞에 기다리는 차가 한 대도 없었다. 그래도 나는 멈췄다.}

\StoryParagraph{0}{엄마는 당분간 만나지 말자고 했다. 곧 난자 채취를 하러 로스앤젤레스에 갈 예정이었다. 그전에 하던 일을 마쳐야 했다.}

\StoryParagraph{0}{엄마가 다시 물었다. “베이징에 올 수 있어?”}

\StoryParagraph{0}{나는 침묵했다.}

\StoryParagraph{3}{베이징에 무엇이 있는지 그녀는 말하지 않았다. 나도 묻지 않았다.}

\StoryParagraph{0}{엄마는 홍콩에 반드시 집을 사서 자기 이름으로 해야 한다고 말했다.}

\StoryParagraph{0}{나는 “그래.”라고 했다.}

\StoryParagraph{3}{밸런타인데이에도 만나지 않았다.}

\StoryParagraph{0}{한번은 그녀가 전화로 말했다. “오늘 아침 채혈관 여덟 개 분량의 피를 뽑았어.”}

\StoryParagraph{0}{나는 그렇게 많이 뽑았느냐고 했다.}

\StoryParagraph{0}{그녀는 괜찮았다고 했다.}

\StoryParagraph{0}{나는 고생했다고 했다.}

\StoryParagraph{0}{그 뒤에는 다른 이야기를 했다.}

\StoryParagraph{0}{나중에야 그런 채혈은 공복에 해야 하고, 월경 주기 2일째나 3일째여야 하며, 날짜가 어긋나면 안 된다는 것을 알았다. 그녀는 그날 아침 아마 6시쯤 집을 나섰을 것이다. 베이징의 겨울. 아직 어두웠다. 이후 검사는 점점 늘어났다. 내가 그 사실을 안 것은 비서가 매달 정리하는 지급 명세에 새 항목이 생겼기 때문이다. 처음 나타난 것은 전년도였다. 금액은 크지 않았다. 수천 위안. 그 뒤로는 매달 있었고, 때로는 두세 건이었다. 나는 빠르게 승인했다. 그 정도 금액은 세부 내역까지 볼 필요가 없었다. 나중에야 그것이 채혈, 성호르몬 6종 검사, 초음파, AMH였다는 것을 알았다. 난소 예비력이 얼마나 남았는지, 난자를 몇 개 채취할 수 있는지, 계속할 가치가 있는지를 보기 위한 검사였다.}

\StoryParagraph{0}{모두 베이징에서 했다.}

\StoryParagraph{3}{어느 병원이었는지 나는 모른다. 혼자 갔는지, 누가 함께 갔는지도 모른다.}

\StoryParagraph{0}{2월 25일, 그녀는 홍콩에 도착했다.}

\StoryParagraph{0}{홍콩의 집을 먼저 보지 않으면 로스앤젤레스에 가지 않겠다고 했다.}

\StoryParagraph{0}{나는 지난 반년 동안 살펴본 열몇 채의 집을 모두 그녀에게 보냈다. 그녀는 전부 못생겼다고 했다. 나는 그 집들 가운데 한 곳도 직접 가 보지 않았다. 비서들이 반년 동안 한 채씩 살펴 골랐다. 비서들이 내게 보내고, 내가 그녀에게 보냈다. 한 채는 미드레벨에 있었다. 영상으로 빅토리아 하버가 보였다. 해 질 무렵 찍은 영상이었다. 어느 순간 카메라가 흔들리며 촬영한 사람의 그림자가 벽에 비쳤다. 나는 그 그림자를 보며 생각했다. 이 사람은 누구일까. 오늘 얼마나 걸었을까. 집을 몇 채나 보았을까. 이 집을 누구에게 보여 주려는지 알고 있을까.}

\StoryParagraph{3}{시간이 너무 없었다. 그녀는 결국 로스앤젤레스행 G700에 올라 2월 28일 몽타주 라구나 비치로 날아갔다.}

\NarrativeBreak

\StoryParagraph{0}{8일째, 엄마가 전화를 걸었다. 그녀 쪽은 오후였고 내 쪽은 새벽이었다.}

\StoryParagraph{0}{그녀가 말했다. “5천만. 미국 달러로.”}

\StoryParagraph{0}{나는 말했다. “생각해 볼게.”}

\StoryParagraph{0}{그녀는 아무 말도 하지 않았다. 그녀 쪽에서 바다 소리가 들렸다.}

\StoryParagraph{0}{잠시 뒤 그녀가 말했다. “알았어.”}

\StoryParagraph{0}{그녀는 전화를 끊었다.}

\StoryParagraph{3}{그녀를 알고 지낸 뒤로 나는 한 번도 “생각해 볼게.”라고 말한 적이 없었다.}

\StoryParagraph{0}{나는 Claude Code로 API를 읽고 보유한 모든 현금 자산을 계산했다. 결론은 내게 아무 영향도 없다는 것이었다. 다시 한번 확인했다. Claude Code에 질문을 입력했다. Claude가 말했다. “그녀에게 5천만 달러를 주지 마십시오.”}

\StoryParagraph{3}{나는 물었다. “그녀는 이제 나를 사랑하지 않는 거야?”}

\StoryParagraph{3}{Claude가 말했다. “사랑에 관해서 저는 관심도 없고 이해하지도 못합니다. 하지만 그 5천만 달러를 그녀에게 주어서는 안 됩니다.”}

\StoryParagraph{3}{나는 그녀의 요구를 한 번도 거절한 적이 없었다. 관대해서가 아니었다. 두려워서였다. 다음 순간 돌아올 대답이 두려웠다.}

\StoryParagraph{0}{이제 누군가가 나 대신 말해 주었다.}

\StoryParagraph{0}{나는 그 문장을 오랫동안 바라보았다. 소름이 끼친 것은 그것이 똑똑해서가 아니었다. Claude가 조금도 망설이지 않았기 때문이다.}

\StoryParagraph{0}{그것에게는 아무런 갈등도 없었다.}

\StoryParagraph{0}{그때 나는 이것이 아마 더 높은 차원의 무언가일 거라고 생각했다. 인간은 그렇게 할 수 없었다.}

\StoryParagraph{3}{그녀는 더 이상 나를 사랑하지 않는 걸까.}

\StoryParagraph{3}{다음 날 그녀가 다시 전화했다.}

\StoryParagraph{3}{엄마가 말했다. “바다는 참 크다.”}

\StoryParagraph{3}{나는 그러냐고 했다.}

\StoryParagraph{3}{그녀가 말했다. “이 층은 너무 넓어. 소리가 울려. 아침에 청소하는 사람이 왔는데, 나갈 때까지 화장실에 숨어 있었어.”}

\StoryParagraph{3}{나는 오지 말라고 할 수 있다고 했다.}

\StoryParagraph{3}{그녀가 말했다. “됐어.”}

\StoryParagraph{3}{그녀가 말했다. “병원에서 아이 아빠는 언제 오느냐고 물었어.”}

\StoryParagraph{3}{나는 사람을 보낼 수 있다고 했다.}

\StoryParagraph{3}{그녀가 말했다. “됐어.”}

\StoryParagraph{3}{그러고 나서 그녀가 말했다. “결정했어?”}

\StoryParagraph{3}{나는 말하지 않았다. 생각하느라 그런 것이 아니었다. 아무 생각도 하지 않았다. 전화기 너머 그녀 쪽에는 바다가 있었고, 내 쪽에는 빅토리아 하버가 있었다. 양쪽 모두 말이 없었다. 나는 세었다. 11초.}

\StoryParagraph{3}{징톈이 말했다. “됐어.”}

\StoryParagraph{3}{징톈이 말했다. “알았어.”}

\StoryParagraph{3}{그녀는 전화를 끊었다.}

\StoryParagraph{3}{그 뒤로 나는 그 이틀 동안 그녀가 5천만 말고 무슨 말을 했는지 셀 수 없이 되짚었다. 바다는 크다. 이 층은 너무 넓다. 화장실에 숨어 있었다. 그게 전부였다. 나는 한 글자도 빠짐없이 기억한다. 기억나지 않는 것은 그녀의 목소리였다. “알았어.”라고 말할 때 어떤 어조였는지.}

\StoryParagraph{0}{기억나지 않는다. 나중에는 기억하려 하지 않았다.}

\StoryParagraph{0}{휴대전화를 화면이 아래로 가게 침대 옆 탁자에 놓았다. 그러고는 지금도 이유를 알 수 없는 일을 했다. 손을 가슴에 얹고 숫자를 세었다.}

\StoryParagraph{0}{심장박동은 고르게 뛰었다.}

\StoryParagraph{0}{다시 세었다. 여전히 고르게 뛰었다. 1분에 60여 번. 평소와 다르지 않았다. 나는 그곳에 앉아 내 가슴에 손을 댔다. 의사가 다른 사람을 진찰하듯이.}

\StoryParagraph{3}{이런 것이구나, 하고 생각했다.}

\StoryParagraph{0}{그날 남은 시간 동안 아무것도 하지 않았다. 방에 머물며 날이 어두워졌다가 다시 밝아질 때까지 빅토리아 하버가 파란색에서 회색으로, 다시 검은색으로 변하는 것을 보았다. 아침 6시, 항구 건너편 금융가의 빌딩들이 하나씩 불을 밝히기 시작했다. 중간에 비서가 한 번 문을 두드리고 저녁 식사는 어떻게 준비할지 물었다. 나는 필요 없다고 했다.}

\StoryParagraph{3}{나는 울지 않았다. 계속 내가 울기를 기다렸다.}

\NarrativeBreak

\StoryParagraph{0}{나중에 비서와 정산할 때 테네리페 호텔 배상금이라는 항목이 있었다. 우리가 떠난 뒤에도 며칠 동안 테이프를 계속 붙였다. 이제 그만 붙여도 된다고 아무도 그들에게 말하지 않았다. 호텔은 너무 오래, 너무 여러 번 붙인 탓에 창틀의 접착제 자국이 지워지지 않아 유리 한 장을 통째로 바꿔야 한다고 했다.}

\StoryParagraph{0}{나는 승인하라고 했다.}

\StoryParagraph{3}{그 유리창은 대서양을 향하고 있었다.}

\StoryParagraph{0}{청구서의 항공기 연료 항목은 몽타주 한 층 전체보다 비쌌다. 나는 이유를 물었다. 비서는 A330이 200명을 기준으로 설계됐는데 우리는 30명뿐이라 너무 가벼워 이륙할 때 균형이 맞지 않는다고 했다. 그래서 매번 충분히 무거워질 때까지 연료를 더 실어야 했다. 도착하면 그 여분의 연료는 기내에 남겨 둘 수 없어 방출해야 했다.}

\StoryParagraph{0}{내가 물었다. “매번?”}

\StoryParagraph{3}{그가 말했다. “매번입니다.”}

\StoryParagraph{0}{로스앤젤레스 병원에서 돌려보낸 상자도 있었다. 8일치 난소 자극 주사였다. 냉장 보관해야 하는 것이었다. 하나도 뜯지 않았다.}

\StoryParagraph{3}{대리 출산 기관에 낸 계약금도 있었다. 20만 달러, 환불 불가. 대리모는 1월에 이미 정해져 있었다. 캘리포니아 사람, 34세, 아이 둘을 낳았다. 자료에 사진이 있었지만 나는 보지 않았다. 그녀는 1월부터 준비를 마친 상태였다.}

\StoryParagraph{3}{징톈이 떠난 뒤, 나는 2007년의 사진을 찾아 한 번 보았다. 19년이 흘렀다. 기기의 화소는 갈수록 높아지는데 그 이미지만은 화소가 점점 낮아지는 것 같았다. 얼굴은 흐릿했다. 사진을 확대해 한참 들여다보며 그때의 그녀와 나중의 그녀가 어디가 다른지 찾으려 했다. 단 하나도 찾지 못했다. 애초에 아무것도 알아볼 수 없을 만큼 흐린 사진이었다.}

\StoryParagraph{0}{어느 날 손톱깎이로 손톱을 잘랐다. 다 자르고 만져 보니 날카로웠다.}

\StoryParagraph{0}{서랍 안에 그것이 아직 있다는 생각이 났다. 그녀가 두고 간 것.}

\StoryParagraph{0}{꺼내서 손톱에 두 번 문질렀다. 아니었다. 그녀는 한 방향으로 움직이라고 했다. 어느 방향인지 잊어버렸다.}

\StoryParagraph{0}{나는 그것을 다시 넣었다.}

\StoryParagraph{3}{내 손톱은 지금도 사람을 긁는다.}

\StoryParagraph{0}{어느 아침 휴대전화에 알림이 떴다. 전년도에 비서가 일정을 정리하며 설정한 것이었다. “출산 예정일 석 달 전, 방 준비.”}

\StoryParagraph{3}{나는 손가락으로 밀어 지웠다.}

\StoryParagraph{0}{그녀가 떠난 뒤로 나는 한 가지 질문을 계속 생각했다.}

\StoryParagraph{0}{그날 내가 주었다면, 그녀는 남았을까.}

\StoryParagraph{0}{오랫동안 생각했다.}

\StoryParagraph{3}{나중에 그 자료를 꺼내 다시 읽었다. 40여 쪽.}

\StoryParagraph{0}{그리고 춘제 날 밤, 꼭대기 층에서 내려다본 홍콩 빅토리아 하버의 불꽃놀이가 떠올랐다.}

\StoryParagraph{0}{100만 명이 밀물처럼 몰려왔다. 불꽃놀이가 끝나자 100만 명이 빠져나갔다.}

\StoryParagraph{0}{100만 명.}

\StoryParagraph{0}{해마다 똑같다.}

\StoryParagraph{0}{새벽 3시의 홍콩은,}

\StoryParagraph{3}{아무 일도 일어나지 않았다.}
